\chapter{Limit dan Kekontinuan}\label{ch:g12:limcont}

Bab ini memperluas gagasan limit dari \omterm{def:g12:seq:sequence}{barisan} ke \omterm{def:g11:func:function}{fungsi} dari satu peubah
real, lalu memperkenalkan \omterm{def:g12:limcont:continuity}{kekontinuan}. Hasil utamanya adalah teorema
nilai antara, yang mengubah keterangan kualitatif (``$f$ \omterm{def:g12:limcont:continuity}{kontinu} dan berganti
tanda'') menjadi keberadaan penyelesaian sebuah \omterm{def:g10:algebra:equation}{persamaan}.

\section{Limit fungsi}

Di sepanjang bab ini, $f$ adalah \omterm{def:g11:func:function}{fungsi} yang terdefinisi pada sebuah \omterm{def:g10:numbers:interval}{interval}
atau gabungan \omterm{def:g10:numbers:interval}{interval} $D \subseteq \R$.

\begin{definition}[Limit di takhingga]\label{def:g12:limcont:liminf}
Misalkan $f$ terdefinisi pada \omterm{def:g10:numbers:interval}{interval} $\intco{a}{+\infty}$ dan $\ell \in \R$.
Kita katakan bahwa $f$ \emph{menuju $\ell$ di $+\infty$}\index{limit!fungsi},
ditulis $\lim\limits_{x\to+\infty} f(x) = \ell$, bila setiap \omterm{def:g10:numbers:interval}{interval} terbuka
yang memuat $\ell$ memuat semua nilai $f(x)$ untuk $x$ yang cukup besar:
untuk setiap $\varepsilon > 0$ ada $A$ sehingga untuk semua $x \geq A$
berlaku $\abs{f(x) - \ell} \leq \varepsilon$.

Kita katakan $f$ menuju $+\infty$ di $+\infty$ bila untuk setiap $M \in \R$ ada
$A$ sehingga $f(x) \geq M$ untuk semua $x \geq A$. Limit di $-\infty$
dan limit yang sama dengan $-\infty$ didefinisikan secara serupa.
\end{definition}

\begin{definition}[Limit di sebuah titik]\label{def:g12:limcont:limpoint}
Misalkan $a \in \R$ dan misalkan $f$ terdefinisi di dekat $a$ (kecuali mungkin
di $a$). Kita katakan bahwa $f$ menuju $\ell \in \R$ di $a$, ditulis
$\lim\limits_{x \to a} f(x) = \ell$, bila untuk setiap $\varepsilon > 0$ ada
$\delta > 0$ sehingga untuk semua $x \in D$ dengan
$\abs{x - a} \leq \delta$ berlaku $\abs{f(x) - \ell} \leq \varepsilon$.

Kita katakan $f$ menuju $+\infty$ di $a$ bila untuk setiap $M$ ada $\delta > 0$
sehingga $f(x) \geq M$ setiap kali $x \in D$ dan $\abs{x - a} \leq \delta$.
Limit sepihak ($x \to a^+$, $x \to a^-$) membatasi $x$ pada $x > a$ atau
$x < a$.
\end{definition}

\begin{example}
$\lim\limits_{x\to+\infty} \frac{1}{x} = 0$,
$\lim\limits_{x\to 0^+} \frac{1}{x} = +\infty$,
$\lim\limits_{x\to 0^-} \frac{1}{x} = -\infty$. \omterm{def:g11:func:function}{Fungsi}
$x \mapsto \frac1x$ tidak berlimit (dua pihak) di $0$.
\end{example}

\begin{definition}[Asimtot]\label{def:g12:limcont:asymptote}
Garis $y = \ell$ disebut \emph{asimtot mendatar}\index{asimtot} kurva
$f$ di $+\infty$ (masing-masing $-\infty$) bila
$\lim_{x\to+\infty} f(x) = \ell$ (masing-masing di $-\infty$). Garis $x = a$
disebut \emph{asimtot tegak}\index{asimtot} bila $f$ mempunyai limit sepihak yang takhingga di $a$.
\end{definition}

\begin{omfigure}
\begin{tikzpicture}
\begin{axis}[omaxis,
  width=8.8cm, height=6cm,
  xmin=-2.9, xmax=4.9, ymin=-5.2, ymax=9.2,
  xtick={1}, ytick={2}]
  \draw[omProp, dashed] (axis cs:-2.9,2) -- (axis cs:4.9,2);
  \draw[omProp, dashed] (axis cs:1,-5.2) -- (axis cs:1,9.2);
  \addplot[omDef, thick, domain=-2.8:0.855, samples=60] {2 + 1/(x-1)};
  \addplot[omDef, thick, domain=1.145:4.8, samples=60] {2 + 1/(x-1)};
\end{axis}
\end{tikzpicture}

{\small Kurva $f(x) = 2 + \frac{1}{x-1}$ beserta \omterm{def:g12:limcont:asymptote}{asimtot mendatarnya}
$y = 2$ (di $\pm\infty$) dan \omterm{def:g12:limcont:asymptote}{asimtot tegaknya} $x = 1$.}
\end{omfigure}

\begin{proposition}[Operasi pada limit]\label{prop:g12:limcont:operations}
Aturan \cref{prop:g12:seq:operations} (jumlah, hasil kali, hasil bagi) berlaku
kata demi kata untuk limit \omterm{def:g11:func:function}{fungsi}, di sebuah titik atau di takhingga, dengan
keempat bentuk taktentu yang sama $(+\infty)+(-\infty)$, $0\times\infty$,
$\frac{\infty}{\infty}$, $\frac{0}{0}$.
\end{proposition}

\begin{proof}
Buktinya sama persis dengan kasus \omterm{def:g12:seq:sequence}{barisan}, dengan mengganti ``untuk $n \geq N$''
oleh ``untuk $x \geq A$'' atau ``untuk $\abs{x-a} \leq \delta$''.
\end{proof}

\begin{theorem}[Limit sebuah komposisi]\label{thm:g12:limcont:composition}
Misalkan $f, g$ dua \omterm{def:g11:func:function}{fungsi} dan $a, b, c$ \omterm{def:g10:numbers:sets}{bilangan real} atau $\pm\infty$. Jika
$\lim\limits_{x \to a} f(x) = b$ dan $\lim\limits_{y \to b} g(y) = c$, maka
\[
\lim_{x \to a} g\bigl(f(x)\bigr) = c .
\]
Pernyataan yang sama berlaku dengan sebuah \omterm{def:g12:seq:sequence}{barisan} sebagai ganti $f$: jika
$u_n \to b$ dan $\lim_{y\to b} g(y) = c$, maka $g(u_n) \to c$.
\end{theorem}

\begin{proof}
Ambil $b, c$ berhingga (kasus lainnya serupa). Misalkan $\varepsilon > 0$.
Ada $\delta > 0$ sehingga $\abs{y - b} \leq \delta$ mengakibatkan
$\abs{g(y) - c} \leq \varepsilon$; dan ada $\delta' > 0$ sehingga
$\abs{x - a} \leq \delta'$ mengakibatkan $\abs{f(x) - b} \leq \delta$.
Merangkai keduanya memberi $\abs{g(f(x)) - c} \leq \varepsilon$ untuk
$\abs{x - a} \leq \delta'$.
\end{proof}

\begin{theorem}[Pembandingan dan apitan]\label{thm:g12:limcont:squeeze}
Misalkan $f, g, h$ \omterm{def:g11:func:function}{fungsi} dan $a \in \R \cup \{\pm\infty\}$.
\begin{enumerate}
  \item Jika $f \leq g$ di dekat $a$ dan $f(x) \to +\infty$ ketika $x \to a$,
        maka $g(x) \to +\infty$.
  \item Jika $f \leq g \leq h$ di dekat $a$ dan $f, h$ menuju limit berhingga
        yang sama $\ell$ di $a$, maka $g(x) \to \ell$ ketika $x \to a$.
\end{enumerate}
\end{theorem}

\begin{proof}
Alasan yang sama seperti untuk \omterm{def:g12:seq:sequence}{barisan} (\cref{thm:g12:seq:squeeze}).
\end{proof}

\begin{method}[Menghitung limit dalam praktik]\label{met:g12:limcont:limits}
\begin{itemize}
  \item Polinomial dan \omterm{def:g11:func:function}{fungsi} rasional di $\pm\infty$: keluarkan pangkat
        tertinggi $x$ sebagai faktor; limitnya adalah limit perbandingan suku
        utamanya.
  \item Bentuk $\frac{0}{0}$ di sebuah titik $a$: keluarkan $(x-a)$ sebagai
        faktor pembilang dan penyebutnya, atau kenali sebuah \omterm{def:g11:deriv:derivative}{turunan}
        $\lim_{x\to a}\frac{f(x)-f(a)}{x-a} = f'(a)$ (\cref{ch:g12:deriv}).
  \item \omterm{def:g11:quad:discriminant}{Akar} kuadrat: kalikan dengan bentuk sekawannya.
  \item Faktor yang berosilasi ($\cos$, $\sin$): pakailah teorema apit.
\end{itemize}
\end{method}

\section{Kekontinuan}

\begin{definition}[Kekontinuan]\label{def:g12:limcont:continuity}
Misalkan $f$ terdefinisi pada \omterm{def:g10:numbers:interval}{interval} $I$ dan $a \in I$. \omterm{def:g11:func:function}{Fungsi} $f$ disebut
\emph{kontinu di $a$}\index{kekontinuan} bila
$\lim\limits_{x \to a} f(x) = f(a)$. \omterm{def:g11:func:function}{Fungsi} itu disebut \emph{kontinu pada $I$}\index{kontinu} bila
kontinu di setiap titik $I$.
\end{definition}

Secara gerak hati, \omterm{def:g10:functions:graph}{grafik} \omterm{def:g11:func:function}{fungsi} yang \omterm{def:g12:limcont:continuity}{kontinu} pada sebuah \omterm{def:g10:numbers:interval}{interval} dapat
digambar tanpa mengangkat pena.

\begin{proposition}[Kekontinuan fungsi yang lazim]\label{prop:g12:limcont:usual}
Polinomial, \omterm{def:g11:func:function}{fungsi} rasional, $x \mapsto \sqrt{x}$, $x \mapsto \abs{x}$,
$\cos$, $\sin$, $\exp$ dan $\ln$ \omterm{def:g12:limcont:continuity}{kontinu} pada setiap \omterm{def:g10:numbers:interval}{interval} yang termuat
dalam \omterm{def:g11:func:function}{daerah asalnya}. Jumlah, hasil kali, hasil bagi (di tempat yang
terdefinisi) dan komposisi \omterm{def:g11:func:function}{fungsi} \omterm{def:g12:limcont:continuity}{kontinu} bersifat \omterm{def:g12:limcont:continuity}{kontinu}.
\end{proposition}

\begin{proof}[Bukti sebagian]
Kestabilannya terhadap jumlah, hasil kali, hasil bagi dan komposisi menyusul
dari teorema limit yang bersesuaian. \omterm{def:g12:limcont:continuity}{Kekontinuan} $x \mapsto x$ dan \omterm{def:g11:func:function}{fungsi}
tetap langsung terlihat dari definisinya, sehingga polinomial dan \omterm{def:g11:func:function}{fungsi}
rasional \omterm{def:g12:limcont:continuity}{kontinu}. \omterm{def:g12:limcont:continuity}{Kekontinuan} \omterm{def:g11:func:function}{fungsi} lazim yang lain dibuktikan pada bab
tempat \omterm{def:g11:func:function}{fungsi} itu disusun, atau diterima tanpa bukti.
\end{proof}

\begin{example}
\emph{\omterm{def:g11:func:function}{Fungsi} lantai} $x \mapsto \floor{x}$ (\omterm{def:g10:numbers:sets}{bilangan bulat} terbesar
$\leq x$) tidak \omterm{def:g12:limcont:continuity}{kontinu} di \omterm{def:g10:numbers:sets}{bilangan bulat} $n$ mana pun: limit kirinya di sana
adalah $n - 1$ sedangkan $\floor{n} = n$.
\end{example}

\begin{omfigure}
\begin{tikzpicture}
\begin{axis}[omaxis,
  width=7.2cm, height=5.6cm,
  xmin=-1.9, xmax=3.5, ymin=-2.5, ymax=2.7,
  xtick={-1,1,2,3}, ytick={-2,-1,1,2}]
  \draw[omDef, thick] (axis cs:-1.5,-2) -- (axis cs:-1,-2);
  \draw[omDef, thick] (axis cs:-1,-1) -- (axis cs:0,-1);
  \draw[omDef, thick] (axis cs:0,0) -- (axis cs:1,0);
  \draw[omDef, thick] (axis cs:1,1) -- (axis cs:2,1);
  \draw[omDef, thick] (axis cs:2,2) -- (axis cs:3,2);
  \fill[omDef] (axis cs:-1,-1) circle (1.5pt);
  \fill[omDef] (axis cs:0,0) circle (1.5pt);
  \fill[omDef] (axis cs:1,1) circle (1.5pt);
  \fill[omDef] (axis cs:2,2) circle (1.5pt);
  \draw[omDef, fill=white] (axis cs:-1,-2) circle (1.5pt);
  \draw[omDef, fill=white] (axis cs:0,-1) circle (1.5pt);
  \draw[omDef, fill=white] (axis cs:1,0) circle (1.5pt);
  \draw[omDef, fill=white] (axis cs:2,1) circle (1.5pt);
  \draw[omDef, fill=white] (axis cs:3,2) circle (1.5pt);
\end{axis}
\end{tikzpicture}

{\small \omterm{def:g11:func:function}{Fungsi} lantai melompat di setiap \omterm{def:g10:numbers:sets}{bilangan bulat}: nilainya (titik penuh)
berbeda dari limit kirinya (titik kosong).}
\end{omfigure}

\begin{proposition}[Barisan dan kekontinuan]\label{prop:g12:limcont:seqcont}
Jika $u_n \to \ell$ dan $f$ \omterm{def:g12:limcont:continuity}{kontinu} di $\ell$, maka $f(u_n) \to f(\ell)$.
Khususnya, jika \omterm{def:g12:seq:sequence}{barisan} yang ditentukan oleh $u_{n+1} = f(u_n)$ \omterm{def:g12:seq:limit}{konvergen} ke
$\ell$ dan $f$ \omterm{def:g12:limcont:continuity}{kontinu} di $\ell$, maka $f(\ell) = \ell$.
\end{proposition}

\begin{proof}
Pernyataan pertamanya adalah \cref{thm:g12:limcont:composition} yang diterapkan
dengan \omterm{def:g12:seq:sequence}{barisan} $(u_n)$. Untuk yang kedua, terapkan limit pada kedua ruas
$u_{n+1} = f(u_n)$: ruas kirinya menuju $\ell$, ruas kanannya menuju
$f(\ell)$, dan limitnya tunggal.
\end{proof}

\section{Teorema nilai antara}

\begin{theorem}[Teorema nilai antara]\label{thm:g12:limcont:ivt}
\index{teorema nilai antara}
Misalkan $f$ \omterm{def:g12:limcont:continuity}{kontinu} pada $\intcc{a}{b}$ dan misalkan $k$ sebarang bilangan di
antara $f(a)$ dan $f(b)$. Maka ada $c \in \intcc{a}{b}$ sehingga
$f(c) = k$.
\end{theorem}

\begin{omfigure}
\begin{tikzpicture}
\begin{axis}[omaxis,
  width=8.8cm, height=6cm,
  xmin=-0.4, xmax=4.4, ymin=-0.5, ymax=4.7,
  xtick={0.5,3.19,3.8}, xticklabels={$a$,$c$,$b$},
  ytick={1.26,2.5,3.98}, yticklabels={$f(a)$,$k$,$f(b)$}]
  \addplot[omProp, dashed, domain=0:3.19] {2.5};
  \addplot[omDef, thick, domain=0.5:3.8] {0.08*x^3 - 0.5*x + 1.5};
  \draw[black, dashed, thin] (axis cs:0.5,0) -- (axis cs:0.5,1.26)
    -- (axis cs:0,1.26);
  \draw[black, dashed, thin] (axis cs:3.8,0) -- (axis cs:3.8,3.98)
    -- (axis cs:0,3.98);
  \draw[black, dashed, thin] (axis cs:3.19,2.5) -- (axis cs:3.19,0);
  \fill (axis cs:0.5,1.26) circle (1.5pt);
  \fill (axis cs:3.8,3.98) circle (1.5pt);
  \fill (axis cs:3.19,2.5) circle (1.5pt);
\end{axis}
\end{tikzpicture}

{\small Teorema nilai antara: kurva \omterm{def:g12:limcont:continuity}{kontinu} yang menghubungkan
$(a, f(a))$ dengan $(b, f(b))$ pasti memotong setiap garis mendatar
$y = k$ yang terletak di antara $f(a)$ dan $f(b)$.}
\end{omfigure}

\begin{proof}
Dengan mengganti $f$ oleh $f - k$ (dan oleh $k - f$ bila perlu), kita boleh
menganggap $f(a) \leq 0 \leq f(b)$ lalu mencari sebuah nol $f$. Kita susun dua
\omterm{def:g12:seq:sequence}{barisan} lewat \emph{dikotomi}\index{dikotomi}: ambil $a_0 = a$, $b_0 = b$;
diberikan $a_n \leq b_n$ dengan $f(a_n) \leq 0 \leq f(b_n)$, misalkan
$m = \frac{a_n + b_n}{2}$ lalu definisikan
\[
(a_{n+1}, b_{n+1}) =
\begin{cases}
(a_n, m) & \text{jika } f(m) \geq 0,\\
(m, b_n) & \text{jika } f(m) < 0.
\end{cases}
\]
Pada kedua kasusnya $f(a_{n+1}) \leq 0 \leq f(b_{n+1})$, \omterm{def:g12:seq:sequence}{barisan} $(a_n)$
naik, $(b_n)$ turun, dan
$b_n - a_n = \frac{b-a}{2^n} \to 0$: jadi kedua \omterm{def:g12:seq:sequence}{barisannya} berdampingan (lihat
\cref{exo:g12:seq:10}) dan \omterm{def:g12:seq:limit}{konvergen} ke limit bersama
$c \in \intcc{a}{b}$. Menurut \omterm{def:g12:limcont:continuity}{kekontinuannya}, $f(a_n) \to f(c)$ dan
$f(b_n) \to f(c)$; menerapkan limit pada $f(a_n) \leq 0$ dan
$f(b_n) \geq 0$ memberi $f(c) \leq 0 \leq f(c)$, sehingga $f(c) = 0$.
\end{proof}

\begin{remark}
Bukti \omterm{thm:g12:limcont:ivt}{dikotominya} adalah sebuah algoritme: bukti itu menghasilkan \omterm{def:g10:numbers:approx}{hampiran}
sebuah penyelesaian setepat yang diinginkan, dengan galat yang terbagi dua pada
setiap langkah. Beginilah komputer dapat menyelesaikan $f(x) = 0$ hanya dengan
mengetahui cara menghitung $f$.
\end{remark}

\begin{theorem}[Teorema bijeksi]\label{thm:g12:limcont:bijection}
Misalkan $f$ \omterm{def:g12:limcont:continuity}{kontinu} dan \emph{tegas \omterm{def:g12:seq:monotonic}{monoton}} pada sebuah \omterm{def:g10:numbers:interval}{interval} $I$. Maka
untuk setiap $k$ yang terletak tegas di antara limit $f$ di titik ujung $I$,
\omterm{def:g10:algebra:equation}{persamaan} $f(x) = k$ mempunyai penyelesaian \emph{tunggal} di $I$.
\end{theorem}

\begin{proof}
Keberadaannya: pilihlah $a, b \in I$ dengan $k$ di antara $f(a)$ dan $f(b)$
(mungkin dilakukan sebab nilai $f$ mendekati limit di titik ujungnya), lalu
terapkan teorema nilai antara pada $\intcc{a}{b}$. Ketunggalannya: jika
$x < y$, kemonotonan yang tegas memberi $f(x) \neq f(y)$, sehingga dua titik
yang berbeda tidak dapat sama-sama menjadi penyelesaian.
\end{proof}

\begin{method}[Menyelesaikan $f(x) = k$ dengan teorema bijeksi]\label{met:g12:limcont:bijection}
Untuk membuktikan bahwa $f(x) = k$ mempunyai tepat satu penyelesaian pada sebuah \omterm{def:g10:numbers:interval}{interval}:
\begin{enumerate}
  \item susunlah \omterm{def:g10:functions:table}{tabel variasi} $f$ (tanda $f'$, lihat
        \cref{ch:g12:deriv});
  \item pada setiap \omterm{def:g10:numbers:interval}{interval} tempat $f$ \omterm{def:g12:limcont:continuity}{kontinu} dan tegas \omterm{def:g12:seq:monotonic}{monoton},
        bandingkan $k$ dengan nilai atau limit di titik ujungnya;
  \item terapkan teorema bijeksi pada setiap \omterm{def:g10:numbers:interval}{interval} semacam itu lalu
        bilanglah penyelesaiannya.
\end{enumerate}
Sebuah kalkulator atau \omterm{thm:g12:limcont:ivt}{dikotomi} lalu menempatkan setiap penyelesaiannya sehalus yang diinginkan.
\end{method}

\begin{example}\label{ex:g12:limcont:cubic}
\omterm{def:g10:algebra:equation}{Persamaan} $x^3 + x = 1$ mempunyai penyelesaian real yang tunggal. Memang
$f(x) = x^3 + x$ \omterm{def:g12:limcont:continuity}{kontinu} dan tegas naik pada $\R$ (sebagai jumlah
\omterm{def:g11:func:function}{fungsi} yang tegas naik), $\lim_{-\infty} f = -\infty$ dan
$\lim_{+\infty} f = +\infty$, sehingga teorema bijeksi berlaku dengan $k = 1$.
Karena $f(0.6) = 0.816 < 1$ dan $f(0.7) = 1.043 > 1$, penyelesaiannya terletak
di $\intoo{0.6}{0.7}$.
\end{example}

\section{Latihan}

\begin{exercise}[$\star$]\label{exo:g12:limcont:1}
Hitunglah limit berikut:
\[
\lim_{x\to+\infty} \frac{2x^3 - x + 1}{5x^3 + x^2}, \qquad
\lim_{x\to-\infty} \bigl(x^2 + 3x - 1\bigr), \qquad
\lim_{x\to 2^+} \frac{x+1}{x-2}, \qquad
\lim_{x\to+\infty} \frac{\cos x}{x}.
\]
\end{exercise}

\begin{exercise}[$\star$]\label{exo:g12:limcont:2}
Misalkan $f(x) = \dfrac{2x^2 - 3x + 1}{x - 1}$, terdefinisi untuk $x \neq 1$.
\begin{enumerate}
  \item Tunjukkan bahwa $f(x) = 2x - 1$ untuk semua $x \neq 1$.
  \item Simpulkan $\lim\limits_{x \to 1} f(x)$. Dapatkah $f$ diperluas menjadi
        \omterm{def:g11:func:function}{fungsi} yang \omterm{def:g12:limcont:continuity}{kontinu} di $1$?
\end{enumerate}
\end{exercise}

\begin{exercise}[$\star$]\label{exo:g12:limcont:3}
Tentukan \omterm{def:g12:limcont:asymptote}{asimtot} kurva
$f(x) = \dfrac{3x - 1}{x + 2}$ lalu sketsalah kurvanya.
\end{exercise}

\begin{exercise}[$\star\star$]\label{exo:g12:limcont:4}
Hitunglah
\[
\lim_{x\to+\infty} \bigl(\sqrt{x^2 + x} - x\bigr)
\qquad\text{dan}\qquad
\lim_{x\to 0} \frac{\sqrt{1+x} - 1}{x}.
\]
\end{exercise}

\begin{exercise}[$\star\star$]\label{exo:g12:limcont:5}
Misalkan $f(x) = x + \sin x$.
\begin{enumerate}
  \item Tunjukkan bahwa $f(x) \to +\infty$ ketika $x \to +\infty$, walaupun
        $f$ tidak \omterm{def:g12:seq:monotonic}{monoton} pada \omterm{def:g10:numbers:interval}{interval} $\intco{A}{+\infty}$ mana pun.
  \item Tunjukkan bahwa $g(x) = \dfrac{x + \sin x}{x - \sin x}$ menuju $1$ di
        $+\infty$.
\end{enumerate}
\end{exercise}

\begin{exercise}[$\star\star$]\label{exo:g12:limcont:6}
Tunjukkan bahwa \omterm{def:g10:algebra:equation}{persamaan} $x^3 - 3x + 1 = 0$ mempunyai tepat tiga penyelesaian
real, lalu berikan untuk masing-masing sebuah \omterm{def:g10:numbers:interval}{interval} sepanjang $1$ yang
memuatnya. (Petunjuk: telaahlah variasi $f(x) = x^3 - 3x + 1$; \omterm{def:g11:deriv:derivative}{turunannya}
adalah $3x^2 - 3$.)
\end{exercise}

\begin{exercise}[$\star\star$]\label{exo:g12:limcont:7}
Misalkan $f \colon \intcc{0}{1} \to \intcc{0}{1}$ \omterm{def:g12:limcont:continuity}{kontinu}. Tunjukkan bahwa $f$
mempunyai \emph{\omterm{pb:g10:functions:1}{titik tetap}}: ada $c \in \intcc{0}{1}$ dengan $f(c) = c$.
(Petunjuk: terapkan teorema nilai antara pada $g(x) = f(x) - x$.)
\end{exercise}

\begin{exercise}[$\star\star\star$]\label{exo:g12:limcont:8}
Seorang pendaki menyusuri jalur gunung, berangkat pukul 08.00 dan tiba pukul
20.00. Keesokan harinya ia menuruni jalur yang sama, lagi-lagi berangkat pukul
08.00 dan tiba pukul 20.00. Tunjukkan bahwa ada saat dalam sehari ketika ia
berada tepat pada titik jalur yang sama pada kedua harinya. (Petunjuk:
perkenalkan selisih kedua \omterm{def:g11:func:function}{fungsi} kedudukannya.)
\end{exercise}

\begin{exercise}[$\star\star\star$]\label{exo:g12:limcont:9}
Misalkan $f(x) = x^5 + x^3 + x - 1$.
\begin{enumerate}
  \item Tunjukkan bahwa $f$ mempunyai satu nol real tunggal $\alpha$, dan
        bahwa $\alpha \in \intoo{0}{1}$.
  \item Perikan sebuah prosedur \omterm{thm:g12:limcont:ivt}{dikotomi} yang menghitung $\alpha$ dengan
        ketelitian $10^{-6}$, lalu berikan banyaknya iterasi yang diperlukan
        bila dimulai dari $\intcc{0}{1}$.
\end{enumerate}
\end{exercise}

\section{Soal: Menyeberangi sungai pasti basah}

\begin{problem}[{Soal akhir pekan --- teorema nilai antara
memaksa titik tetap, antipoda bersuhu sama dan meja yang berhenti
goyang: tiga teorema yang terdengar mustahil dari satu gagasan}]
\label{pb:g12:limcont:1}
Remaslah sebuah \omterm{def:g10:functions:function}{peta} lalu jatuhkan di mana saja di atas wilayah yang
digambarnya: ada satu titik \omterm{def:g10:functions:function}{peta} yang terletak \emph{tepat} di atas tempat
yang diwakilinya. Pada setiap saat, dua titik khatulistiwa yang berhadapan
secara diametral mempunyai suhu yang \emph{persis} sama. Dan sebuah meja
persegi yang goyang di atas tanah bergelombang selalu dapat ditenangkan dengan
\emph{memutarnya}. Tiga sulap pesta, satu teorema: besaran yang \omterm{def:g12:limcont:continuity}{kontinu}
tidak dapat beralih dari negatif ke positif tanpa lenyap
(\cref{thm:g12:limcont:ivt}). Soal ini melatih limit, lalu
membangun ketiga keajaiban itu --- sekaligus mengetahui persis apa yang
tidak dijanjikan teoremanya.

\medskip
\noindent\textbf{Bagian I --- Kefasihan dengan limit.}
\begin{enumerate}
  \item Hitunglah $\lim_{x \to +\infty} \dfrac{3x^2 - x}{x^2 + 5}$,
        \ $\lim_{x \to +\infty} \dfrac{2x + 1}{x^2 + 1}$, \ dan
        $\lim_{x \to +\infty} (x^3 - x^2)$
        (\cref{met:g12:limcont:limits}).
  \item Berikan \omterm{def:g12:limcont:asymptote}{asimtot}
        $f(x) = \dfrac{3x - 1}{x + 2}$
        (\cref{def:g12:limcont:asymptote}), dengan menghitung
        limit yang bersangkutan.
  \item Hitunglah $\lim_{x \to 3} \dfrac{x^2 - 9}{x - 3}$
        (dengan \omterm{def:g10:algebra:expand}{memfaktorkan}) dan
        $\lim_{x \to 0} \dfrac{\sqrt{x + 1} - 1}{x}$
        (dengan bentuk sekawan).
  \item Dengan teorema apit
        (\cref{thm:g12:limcont:squeeze}): hitunglah
        $\lim_{x \to 0} x \sin\frac1x$ dan
        $\lim_{x \to +\infty} \frac{\sin x}{x}$.
  \item Lewat komposisi dan suku yang menguasai:
        $\lim_{x \to +\infty} \dfrac{\sqrt{4x^2 + x}}{x}$.
\end{enumerate}

\medskip
\noindent\textbf{Bagian II --- \omterm{def:g11:quad:discriminant}{Akar}, yang tersahkan.}
\begin{enumerate}[resume]
  \item Tunjukkan bahwa \omterm{def:g10:algebra:equation}{persamaan} $x^3 + x = 5$ mempunyai paling sedikit satu
        penyelesaian di $\intcc{1}{2}$.
  \item Tunjukkan bahwa penyelesaiannya tunggal pada $\R$
        (\cref{thm:g12:limcont:bijection}).
  \item Buktikan yang klasik: \emph{setiap polinomial berderajat ganjil
        mempunyai paling sedikit satu \omterm{def:g11:quad:discriminant}{akar} real}. (Tuliskan alasannya
        untuk $p(x) = x^3 + a x^2 + b x + c$: hitunglah kedua
        limit takhingganya lewat suku yang menguasai, lalu terapkan teorema
        nilai antara pada \omterm{def:g10:numbers:interval}{interval} yang besar.)
  \item Menempatkan \omterm{def:g11:quad:discriminant}{akar} pertanyaan 6 lewat \omterm{thm:g12:limcont:ivt}{dikotomi}: berapa kali
        $\intcc{1}{2}$ harus dibagi dua agar galatnya terjamin di bawah
        $10^{-6}$ (\cref{exo:g12:limcont:9})? Bandingkan dengan
        penggandaan angka pada metode Newton
        (\cref{pb:g11:deriv:1}): kira-kira berapa langkah yang
        diperlukannya?
  \item Contoh penyangkal yang menjaga teoremanya:
        $g(x) = \frac1x$ memenuhi $g(-1) = -1 < 0$ dan
        $g(1) = 1 > 0$, namun tidak pernah lenyap. Hipotesis mana pada
        teorema nilai antara yang gagal --- dan apa yang diajarkan contoh itu
        tentang memeriksa hipotesis?
\end{enumerate}

\medskip
\noindent\textbf{Bagian III --- Tiga teorema yang terdengar
mustahil.}
\begin{enumerate}[resume]
  \item Teorema \omterm{pb:g10:functions:1}{titik tetap} pada sebuah ruas: jika
        $f : \intcc{0}{1} \to \intcc{0}{1}$ \omterm{def:g12:limcont:continuity}{kontinu}, buktikan
        bahwa ada $c$ yang memenuhi $f(c) = c$. (Telaahlah
        $g(x) = f(x) - x$ di kedua ujungnya.)
  \item Terjemahkan pertanyaan 11 ke \omterm{def:g10:functions:function}{peta} yang teremas: apa yang berperan
        sebagai $f$, mengapa nilainya mendarat di
        $\intcc{0}{1}$ (secara kiasan: di dalam wilayahnya), dan apa
        titik yang terletak tepat di atas letaknya sendiri? (Satu
        dimensi diterima diam-diam --- pernyataan dua dimensi yang jujur
        adalah teorema Brouwer, pada jilid universitas.)
  \item Khatulistiwa: misalkan $T$ sebuah suhu yang \omterm{def:g12:limcont:continuity}{kontinu} pada sebuah
        lingkaran, yang dipandang sebagai \omterm{def:g11:func:function}{fungsi} berkala $2\pi$ dengan
        $T(0) = T(2\pi)$. Ambil $g(t) = T(t) - T(t + \pi)$.
        Hitunglah $g(0) + g(\pi)$, simpulkan bahwa $g$ lenyap
        di suatu tempat pada $\intcc{0}{\pi}$, lalu nyatakan teoremanya
        tentang titik antipoda.
  \item Meja yang goyang: sebuah meja persegi yang sempurna kaku berdiri
        di atas tanah bergelombang yang \omterm{def:g12:limcont:continuity}{kontinu} (tanpa undakan!); tiga kakinya
        selalu menyentuh tanah, satu melayang setinggi $h(\theta)$ yang
        bergantung secara \omterm{def:g12:limcont:continuity}{kontinu} pada sudut putar meja itu,
        $\theta$. Jelaskan mengapa memutarnya seperempat putaran memaksa
        celah yang melayang itu berganti tanda --- lalu simpulkan bahwa
        ada sudut $\theta^*$ yang menenangkan keempat kakinya. (Alasan ini
        diterbitkan sebagai teorema sungguhan pada tahun 2005;
        pemilik kafe sudah lebih dahulu tahu kiatnya.)
  \item Pendaki pada \cref{exo:g12:limcont:8} memakai skema yang
        sama. Nyatakan pola yang dibagi pertanyaan 11, 13,
        14 dan pendaki itu --- bangunlah \omterm{def:g11:func:function}{fungsi} yang mana, periksalah apa
        di kedua ujungnya, panggillah teorema apa --- dalam tiga baris.
  \item Apa yang \emph{tidak} diberikan teorema nilai antara: keberadaannya
        datang cuma-cuma, tetapi dua pertanyaan mana yang dibiarkannya
        terbuka, dan perkakas mana pada soal ini yang menjawabnya
        (pertanyaan 7 dan 9)?
\end{enumerate}

\medskip
\noindent\textbf{Bagian IV --- Potret beserta \omterm{def:g12:limcont:asymptote}{asimtotnya}.}
\begin{enumerate}[resume]
  \item Tulislah ulang $f(x) = \dfrac{3x - 1}{x + 2}$ sebagai
        $3 - \dfrac{7}{x + 2}$ lalu pakailah untuk memerikan kurvanya
        selengkapnya: \omterm{def:g12:limcont:asymptote}{asimtot}, variasi pada masing-masing sisi, dan
        sketsanya. (Sebuah \omterm{def:g11:func:reference}{hiperbola} yang tergeser --- kurva biaya rata-rata
        pada \cref{pb:g10:reffunc:1} pun demikian.)
  \item \omterm{def:g11:func:function}{Fungsi} $f(x) = \dfrac{x^2 - 1}{x - 1}$ tidak
        terdefinisi di $1$. Nilai mana yang diberikan kepada $f(1)$ yang
        membuatnya \omterm{def:g12:limcont:continuity}{kontinu} di sana --- dan mengapa para analis menyebut titik
        semacam itu diskontinuitas yang \emph{terhapuskan}?
  \item \omterm{def:g11:func:function}{Fungsi} lantai $x \mapsto \lfloor x \rfloor$
        (\omterm{def:g10:numbers:sets}{bilangan bulat} terbesar $\leq x$): di mana \omterm{def:g11:func:function}{fungsi} itu \omterm{def:g12:limcont:continuity}{kontinu},
        di mana ia melompat, dan mengapa kesimpulan teorema nilai antara
        gagal untuknya (carilah $a < b$ dengan
        $\lfloor a \rfloor < \frac12 < \lfloor b \rfloor$ namun
        tanpa penyelesaian $\lfloor x \rfloor = \frac12$)?
  \item Penutup --- kedua wajah \omterm{def:g12:limcont:continuity}{kekontinuan}: janji lokalnya
        (nilai $=$ limit: tanpa teleportasi) dan kekuatan globalnya
        (teorema nilai antara: setiap nilai antaranya disinggahi). Golongkan
        keempat keajaiban Bagian III di bawah wajah yang tepat, lalu berikan
        nama matematis yang tepat bagi semboyan penyeberangan sungai itu.

\end{enumerate}
\end{problem}
